Copula giúp hiểu về sự phụ thuộc ở mức độ sâu hơn, cho phép chúng ta thấy những cạm bẫy tiềm ẩn của các phương pháp tiếp cận sự phụ thuộc vốn chỉ tập trung vào tương quan, đồng thời chỉ cho chúng ta cách xác định một số thước đo hữu ích. Copula tạo điều kiện cho phương pháp tiếp cận mô hình đi từ dưới lên. Cách tiếp cận Copula cho phép chúng ta kết hợp các mô hình biên đã phát triển kỹ lưỡng với nhiều mô hình phụ thuộc khả thi khác nhau, từ đó nghiên cứu mức độ nhạy cảm của rủi ro đối với các đặc điểm phụ thuộc được chỉ định. Vì các loại Copula rất dễ mô phỏng nên chúng đặc biệt phù hợp cho các nghiên cứu rủi ro bằng phương pháp Monte Carlo.

\paragraph{Các thuộc tính cơ bản:}
Một Copula $d$-chiều là một hàm phân phối trên khối lập phương đơn vị $[0, 1]^d$ với các phân phối biên đều tiêu chuẩn. Gọi $C$ là một ánh xạ có dạng $C: [0, 1]^d \to [0, 1]$. Để $C$ là một Copula thì nó phải thỏa mãn ba điều kiện: đầu tiên là $C(u_1, \dots, u_d)$ tăng trong mỗi thành phần $u_i$; tiếp theo là $C(1, \dots, 1, u_i, 1, \dots, 1) = u_i$ với mọi $i \in \{1, \dots, d\}$ và $u_i \in [0, 1]$; cuối cùng là phải thỏa mãn bất đẳng thức hình chữ nhật với mọi $(a_1, \dots, a_d)$ và $(b_1, \dots, b_d) \in [0, 1]^d$ thỏa $a_i \le b_i$, ta có:
\bea
    \sum_{i_1=1}^{2} \dots \sum_{i_d=1}^{2} (-1)^{i_1 + \dots + i_d} C(u_{1i_1}, \dots, u_{di_d}) \ge 0
\eea
Trong đó $u_{j1} = a_j$ và $u_{j2} = b_j$ với mọi $j \in \{1, \dots, d\}$.\\
Đối với thuộc tính biên, với bất kỳ số nguyên $k$ nào thỏa mãn $2 \le k < d$, các phân phối biên $k$-chiều của một Copula $d$-chiều thì bản thân chúng cũng là các Copula. Nếu ta gọi $G$ là một hàm phân phối và $G^{-1}$ là nghịch đảo tổng quát của $G$, tức là hàm $G^{-1}(y) = \inf \{x: G(x) \ge y\}$, thì ta có các biến đổi thành phần. Đầu tiên là biến đổi phân vị, nếu $U \sim U(0, 1)$ có phân phối đều tiêu chuẩn thì $P(G^{-1}(U) \le x) = G(x)$. Tiếp theo là biến đổi xác suất, nếu $Y$ có phân phối tích lũy $G$, trong đó $G$ là một hàm phân phối đơn biến liên tục, thì $G(Y) \sim U(0, 1)$.
\paragraph{Định lý Sklar:}
Định lý này được Sklar giới thiệu vào năm 1959, là định lý quan trọng nhất trong nghiên cứu về Copula và đóng vai trò nền tảng trong việc xây dựng các hàm phân phối đa biến \cite{sklar1959fonctions, nelsen2006introduction}. Định lý thiết lập mối quan hệ giữa hàm phân phối đồng thời $F$ và các hàm phân phối biên $F_1, \dots, F_d$ thông qua một hàm Copula $C$. Về sự tồn tại, với bất kỳ hàm phân phối đồng thời $F$ có biên $F_1, \dots, F_d$ luôn tồn tại một Copula $C: [0, 1]^d \to [0, 1]$ sao cho với mọi $x_1, \dots, x_d$ trong $\bar{R}$, ta có:
\bea
F(x_1, \dots, x_d) = C(F_1(x_1), \dots, F_d(x_d))
\eea
Về tính duy nhất, nếu các biên $F_1, \dots, F_d$ là liên tục thì Copula $C$ là duy nhất. Nếu các biên không liên tục (ví dụ phân phối rời rạc), Copula $C$ vẫn tồn tại nhưng chỉ được xác định duy nhất trên phạm vi giá trị của các biên $(\text{Ran}F_1 \times \dots \times \text{Ran}F_d)$, với $\text{Ran}F_i = F_i(\bar{R})$ biểu thị phạm vi giá trị mà hàm phân phối biên $F_i$ có thể nhận được. Chiều ngược lại mang tính hướng tổng quát, nếu $C$ là một Copula và $F_1, \dots, F_d$ là các hàm phân phối đơn biến, thì hàm $F$ được xác định theo công thức trên chắc chắn là một hàm phân phối đồng thời hợp lệ với các biên chính là $F_1, \dots, F_d$. Mệnh đề tiếp theo chỉ ra rằng nếu $T$ liên tục thì $T \circ T^{-1}(y) = y$. Từ một hàm phân phối đồng thời $F$ có các biên liên tục, bằng cách sử dụng các hàm ngược $F_i^{-1}$ tại $x_i = F_i^{-1}(u_i)$, với $0 \le u_i \le 1$, $i = 1, \dots, d$, ta thu được công thức:
\bea
C(u_1, \dots, u_d) = F(F_1^{-1}(u_1), \dots, F_d^{-1}(u_d))
\eea
Công thức này cho thấy Copula diễn tả sự phụ thuộc trên thang đo phân vị.\\
Ý nghĩa của định lý trong quản trị rủi ro là rất lớn. Định lý cho phép nhà quản lý rủi ro tách biệt việc mô hình hóa hành vi biên của từng yếu tố rủi ro riêng lẻ khỏi việc mô hình hóa cấu trúc phụ thuộc giữa chúng. Sự linh hoạt này giúp nhà nghiên cứu có thể tự do kết hợp các loại phân phối biên khác nhau (như normal, t-Student, Pareto) với các cấu trúc phụ thuộc khác nhau nhằm tạo ra các mô hình phân phối đa biến phức tạp và thực tế hơn. Ngoài ra, khi ứng dụng trong các kịch bản cực đoan, Copula giúp mô hình hóa sự phụ thuộc của các kết quả cực đoan (như trong tính toán Value-at-Risk), điều mà các thước đo tương quan tuyến tính truyền thống thực hiện không tốt. Khi vectơ ngẫu nhiên $X$ có phân phối tích lũy đồng thời $F$ với các phân phối biên liên tục $F_1, \dots, F_d$, thì Copula của $F$ (hoặc $X$) là hàm phân phối đa biến $C$ của $(F_1(X_1), \dots, F_d(X_d))$.

\paragraph{Tính bất biến:}
Copula của một phân phối sẽ không thay đổi dưới các phép biến đổi tăng nghiêm ngặt của các biến biên. Điều này có nghĩa là cấu trúc phụ thuộc (được biểu diễn bởi Copula) tách biệt hoàn toàn với hình dạng của các phân phối biên riêng lẻ. Chúng ta có thể thay đổi thang đo của từng tài sản mà không làm thay đổi bản chất mối quan hệ phụ thuộc giữa chúng. Mệnh đề cụ thể phát biểu rằng: Gọi $(X_1, \dots, X_d)$ là một vectơ ngẫu nhiên với các biên liên tục và Copula $C$, gọi $T_1, \dots, T_d$ là các hàm tăng ngặt. Khi đó $(T_1(X_1), \dots, T_d(X_d))$ cũng có Copula là $C$. \\
Để chứng minh, đầu tiên biến đổi $T_i(X_i)$ có hàm phân phối liên tục $\bar{F}_i(y) := F_i \circ T_i^{-1}(y)$. Để thấy điều này, áp dụng mệnh đề $T \circ T^{-1}(y) \ge y$ ta có:
$$\bar{F}_i(y) = P(X_i \le T_i^{-1}(y)) = P(T_i^{-1} \circ T_i(X_i) \le T_i^{-1}(y))$$
Vì $T_i^{-1}$ là một phép biến đổi tăng, chúng ta sử dụng bổ đề: Nếu $F$ là hàm phân phối tích lũy của biến ngẫu nhiên $X$, thì $P(F(X) \le F(x)) = P(X \le x)$. Ta được:
$$\bar{F}_i(y) = P(T_i(X_i) \le y) + P(X_i = T_i^{-1}(y), T_i(X_i) > y)$$
Nhưng xác suất thứ hai ở vế phải bằng không vì $F_i$ liên tục. Vì $C$ là Copula của $X$, nên có thể tính toán rằng:
$$C(u_1, \dots, u_d) = P(F_1(X_1) \le u_1, \dots, F_d(X_d) \le u_d) = P(\bar{F}_1(T_1(X_1)) \le u_1, \dots, \bar{F}_d(T_d(X_d)) \le u_d)$$
Vì $\bar{F}_i \circ T_i(x) = F_i \circ T_i^{-1} \circ T_i(x) = F_i(x)$ theo mệnh đề đã nêu, suy ra $C$ cũng là Copula của $(T_1(X_1), \dots, T_d(X_d))$.

\paragraph{Giới hạn Fréchet:}
Đối với mọi Copula $C(u_1, \dots, u_d)$, chúng ta luôn có các giới hạn biên sau:
\bea
\max\left\{ \sum_{i=1}^{d} u_i + 1 - d, \, 0 \right\} \le C(u) \le \min\{u_1, \dots, u_d\}
\eea
Mặc dù đưa ra các giới hạn Fréchet cho một Copula, các giới hạn Fréchet có thể đưa ra cho bất kỳ hàm phân phối đa biến nào. Đối với một hàm phân phối tích lũy đa biến $F$ với các biên $F_1, \dots, F_d$, chúng ta thiết lập bằng lập luận tương tự là:
\bea
\max\left\{ \sum_{i=1}^{d} F_i(x_i) + 1 - d, \, 0 \right\} \le F(x) \le \min\{F_1(X_1), \dots, F_d(X_d)\}
\eea
Vì vậy chúng ta có các giới hạn cho $F$ theo các phân phối biên riêng lẻ của nó.
